\par\indent\SourceLeaf{236}%
III. Setting \(x=r\cos\varphi,\ y=r\sin\varphi\) gives
\[
z=\tfrac{1}{2}(T\cos\varphi^2+V\sin\varphi^2)rr+\Omega
\]
It follows that a section by a plane normal to the surface at \(A\) and inclined to the \(x\)-axis at an angle \(\varphi\) gives a plane curve whose radius of curvature at \(A\) is
\[
=\frac{1}{T\cos\varphi^2+V\sin\varphi^2}
\]
\par\indent IV. Thus, whenever \(T=V\), the radii of curvature in \textit{all} normal planes will be equal. If \(T\) and \(V\) are unequal, then, since \(T\cos\varphi^2+V\sin\varphi^2\), for any value of the angle \(\varphi\), lies between \(T\) and \(V\), the radii of curvature in the principal sections considered in I and II evidently correspond to the extreme curvatures. If \(T\) and \(V\) have the same sign, one corresponds to the greatest curvature and the other to the least. If \(T\) and \(V\) have opposite signs, one corresponds to the greatest convexity and the other to the greatest concavity. These conclusions contain almost all that the distinguished \textsc{Euler} first taught about the curvature of curved surfaces.

V. The measure of curvature of the curved surface at \(A\) takes the very simple form \(k=TV\). This gives us the following theorem.

\textsc{Theorem.} \textit{The measure of curvature at any point of the surface equals a fraction with numerator one and denominator the product of the two extreme radii of curvature in sections by normal planes.}

It is also clear that the measure of curvature is positive for concave-concave or convex-convex surfaces (the distinction is inessential) and negative for concave-convex surfaces. If the surface consists of parts of both kinds, the measure of curvature must vanish at the boundaries between them. We shall discuss in more detail below the nature of curved surfaces on which the measure of curvature vanishes everywhere.

\Article{9}
The general formula for the measure of curvature given at the end of Article 7 is the simplest of all, since it involves only five elements. We arrive at a more complicated formula, involving nine elements, if we choose to use
